\chapimg{fig/at/opq_at22.jpg}
\chapterlead{\subhd{After this chapter you can} say how the study behind SWAFL was built and checked: what each of the six papers asked, which equation answers it, which evidence supports the answer, and how the same chain of ideas runs from the memory of soil water through the audit of the model to the decision tools of the application.}
\chapter{The Research Behind SWAFL: Evidence, Derivation and Validation}\label{chap:22}
\markboth{The Research Behind SWAFL}{The Research Behind SWAFL}
\IfFileExists{extra/th22.tex}{\section{The structure of a research paper}
A scientific paper reports one study in a fixed order: the \emph{abstract} (the whole in 200 words), the \emph{introduction} (the question and why it matters), the \emph{methods} (what was done, in enough detail to repeat it), the \emph{results} (what was found, with the figures and tables), the \emph{discussion} (what it means and where it fails), the \emph{conclusions} and the \emph{references}. The figures carry the results: a good figure can be understood with its caption alone, and the caption says what is plotted, what the panels (a), (b) and (c) are, and what the symbols and colours mean.

\section{How the six papers fit together}
The papers report one project from six sides. The first states the question (what happens to the water stored in the soil when it is released) and answers with the rate-competition law. The second audits the model (does the balance close, what is the term that is not attributed) and gives a compact law of evapotranspiration. The third presents the software that applies the results (the engines and their verification). The fourth asks about the sensor and the trigger (the reach, the fate of a millimetre, the dose). The fifth is about the workbench and the hand-over from a simulation to a decision, and the sixth is about the fields tools and the plants. Read in this order they go from physics to model to tool to use.

\section{Provenance and reproducibility}
A result is \kw{reproducible} if another person, given the data and the code, obtains the same numbers. For this reason every figure of the study is made by a script, every script reads the table of runs, and the table is archived (Zenodo, with a persistent identifier). The \kw{provenance} of a number is the chain from the figure back to the equation, the script, the table and the Daisy run that made it, which the catalogue of figures records. A figure that cannot be traced in this way is an illustration, not a result.

\section{Highlights and graphical abstracts}
A \kw{highlight} is a single sentence of a main finding; a set of three to five gives the reader the content in a minute. A \kw{graphical abstract} is a single picture that shows the question and the answer without text; it is the first thing most readers see, and its aim is to say what the paper is about, not to give the proof. The graphical abstracts of the papers are collected at the end of this chapter.
}{}
The study behind SWAFL is reported in six papers that look at the same ensemble of Daisy runs from six sides. This chapter follows them in order. Each section states the question of the paper, the equation that answers it, the evidence on which the answer rests and the way the result enters the application; the figures that carry the evidence follow, each with its own caption. A reader who has worked through Chapters~\ref{chap:01} to~\ref{chap:17} will find that every figure shows an equation or a process that the book has already derived; the title of each figure names the idea it shows.
\section{The memory of soil water and the rate-competition law}
\textbf{The question.} A soil that starts a season wet releases some of its stored water during the season. Where does that water go: to transpiration, to evaporation, or below the roots as drainage? The first study measured the answer on a designed ensemble of Daisy runs for ten soils, with the initial suction and the rainfall varied systematically.

\par\noindent\textbf{The answer.} The drained share, the \emph{memory fate} $\phi_D$, runs from about $0.85$ in a sand to $0.02$ in a silty clay loam. It follows the pore-size parameter $n$ (Spearman $\rho=0.95$) and the saturated conductivity ($0.97$), and it is described by a rate-competition law, $\phi_D=K_s/(K_s+q^{\ast})$, in which the soil's ability to pass water competes with the rate at which the water arrives. The law predicts soils that were left out of the fit with an error of $0.058$ (leave-one-soil-out), and it transfers to independent experiments on evaporation and mulch ($\rho=0.85$ and $0.89$).

\par\noindent\textbf{How to read the figures.} They go from the framework, through the hydraulic properties of the ten soils and the grid of storage release, to the fates, the predictors, the law, its modifiers and its robustness. The soils are always ordered by $n$, so that the eye can follow the trend from sand to clay.

\par\vspace{6pt}\noindent\hd{Conceptual framework of the memory of soil water}\par\nobreak
\fig{at/a22_01.jpg}{0.84}{Conceptual and analytical framework. Antecedent storage released from the root zone can leave as drainage, soil evaporation, transpiration or other evaporation (memory fates $\phi$, estimated from E9). An imposed change in evaporative loss is offset by drainage, irrigation or transpiration (loss fates $\psi$, estimated from E16 and E7). The recoverable fraction of an evaporation saving is $1-\psi_D$.}\label{fig:a22_1}
\par\vspace{6pt}\noindent\hd{Hydraulic properties of the ten soils}\par\nobreak
\fig{at/a22_02.jpg}{0.84}{Hydraulic properties of the ten soils from the equation of the paper. The ordering by shape parameter $n$ is the ordering on which the main results are organised.}\label{fig:a22_2}
\par\vspace{6pt}\noindent\hd{The memory grid of storage release}\par\nobreak
\fig{at/a22_03.jpg}{0.84}{Storage release $\bar R$ (mm\,yr$^{-1}$) versus initial suction $h_0$ in the E9 lattice, one panel per soil (a--j, ordered by decreasing van Genuchten $n$). Lines are rain scales $P=0.5$--1.4 (colour bar). Wetter starts (small $h_0$) release more storage; the response saturates early in coarse soils and extends over the full $h_0$ range in fine soils.}\label{fig:a22_3}
\par\vspace{6pt}\noindent\hd{Memory fates and loss fates}\par\nobreak
\fig{at/a22_04.jpg}{0.84}{(a) Memory fates $\phi$ of released antecedent storage (E9; error bars: bootstrap 95\,\% CI for drainage). (b) Loss fates $\psi_D$ (drainage) and $\psi_I$ (irrigation) of an evaporation-factor increase (E16). (c) Drainage share of the evaporation saving from mulch (E7). Soils are ordered by decreasing $\phi_D$.}\label{fig:a22_4}
\par\vspace{6pt}\noindent\hd{Predictors of the memory fate}\par\nobreak
\fig{at/a22_05.jpg}{0.84}{Relationship between the memory fate $\phi_D$ (E9) and four soil-level predictors: (a) van Genuchten $n$, (b) $\log_{10}K_s$ (eight soils), (c) $\TAW$ to 1\,m, (d) effective saturation at 100\,cm. Spearman $\rho$ given in each panel.}\label{fig:a22_5}
\par\vspace{6pt}\noindent\hd{The rate-competition law}\par\nobreak
\fig{at/a22_06.jpg}{0.84}{Rate-competition law. (a) $\phi_D$ against $K_s$ with the fitted curve $K_s/(K_s+q^\ast)$; crosses mark implied $K_s$ for S01 and S02. (b) Leave-one-soil-out predictions against observations (RMSE 0.058). (c) LOSO RMSE of each predictor of $\phi_D$ (asterisk: eight soils; the other predictors are shown for all ten soils, and like-for-like eight-soil values are given in the equation of the paper).}\label{fig:a22_6}
\par\vspace{6pt}\noindent\hd{Modifiers of the drainage share}\par\nobreak
\fig{at/a22_07.jpg}{0.84}{Modifiers of the drainage share $\psi_D$ in E16. Each line is a soil (colour: $n$); (a) rainfall scale, (b) warming, (c) trigger suction, (d) irrigation dose.}\label{fig:a22_7}
\par\vspace{6pt}\noindent\hd{Cross-experiment transfer}\par\nobreak
\fig{at/a22_08.jpg}{0.84}{Cross-experiment transfer. (a) $\psi_D$ from E16 against $\phi_D$ from E9; (b) $\psi_D$ from the mulch lattice E7 against $\phi_D$; (c) E7 against E16. Each point is a soil, coloured by $n$ (as in the equation of the paper).}\label{fig:a22_8}
\par\vspace{6pt}\noindent\hd{The recoverable fraction}\par\nobreak
\fig{at/a22_09.jpg}{0.84}{Recoverable fraction $1-\psi_D$ by soil for the E16 and E7 contrasts.}\label{fig:a22_9}
\par\vspace{6pt}\noindent\hd{Robustness of the key associations}\par\nobreak
\fig{at/a22_10.jpg}{0.84}{Robustness of the key associations. (a) Permutation distribution of Spearman $\rho$ between $n$ and $\phi_D$ (solid line: observed; dotted: its mirror). (b) Bootstrap 95\,\% intervals of five rank associations (10\,000 resamples).}\label{fig:a22_10}
\par\vspace{6pt}\noindent\hd{The climate surface of the drainage share}\par\nobreak
\fig{at/a22_11.jpg}{0.84}{Climate surface of the drainage share $\psi_D$ for coarse loam (S01) from the E11 lattice. (a) $\psi_D$ over rain scale and warming (white contours), (b) $\psi_D$ against $P$ for selected warming levels, (c) $\psi_D$ against warming for selected rain scales.}\label{fig:a22_11}
\par\vspace{6pt}\noindent\hd{Drivers of the soil hydraulic properties}\par\nobreak
\fig{at/a22_12.jpg}{0.84}{Drivers of soil hydraulic properties and their coupling to the soil--plant--atmosphere continuum. (a) Memory fate $\phi_D$ as a function of saturated conductivity $K_s$: simulated values for the eight soils with literature $K_s$ (points) and the rate-competition law (line); arrows show the implied change for a threefold increase or decrease in $K_s$ from 20\,cm\,d$^{-1}$ (law implication, not simulation).}\label{fig:a22_12}
\par\vspace{6pt}\noindent\hd{Structure over time and with management}\par\nobreak
\fig{at/a22_13.jpg}{0.84}{Implied change of the memory fate with soil structure over time and with management, from the fitted rate-competition law (the equation of the paper, $q^\ast=71.5$\,cm\,d$^{-1}$). These are illustrative scenarios under assumed multipliers and rates, not simulation results.}\label{fig:a22_13}
\par\vspace{6pt}\noindent\hd{Evidence on structure, time and mulching}\par\nobreak
\fig{at/a22_14.jpg}{0.84}{Simulation evidence on structure, time and mulching (Daisy ensemble, mean annual values). (a) Run set E7, strong against weak mulch (vapour-flux factor 0.3 against 0.9, capacity 2\,mm): evaporation saved, irrigation saved and extra drainage by soil. (b) Run set E5, rainfed drainage of four finer soils against the consolidation rate $\tau_c$ for conventional tillage (solid) and the no-till value (dotted).}\label{fig:a22_14}
\par\vspace{6pt}\noindent\hd{Check of the van Genuchten parameters in one run}\par\nobreak
\fig{at/a22_15.jpg}{0.84}{One Daisy 7.1.14 run of the loam written by SWAFL 1.1.0 with the van Genuchten--Mualem parameters of every horizon as given (blue) and scaled toward a faster-draining soil (orange: $K_\mathrm{sat}\times4$, $\alpha\times2$, $(n-1)\times1.3$, $\theta_s\times1.1$, $l+1$). (a) Cumulative matrix percolation, (b) cumulative actual evapotranspiration, (c) totals of the two seasons with the difference. Daily logs of Daisy, unchanged; script \mbox{reproducibility/fig\_vg\_check.py}.}\label{fig:a22_15}
\exbox{The drainage share from the rate-competition law}{In Paper 1 the memory fate of the released storage is $\phi_D=K_s/(K_s+q^\ast)$ with $q^\ast=71.5$ cm\,d$^{-1}$. Find $\phi_D$ for a soil with $K_s=300$ cm\,d$^{-1}$ and for one with $K_s=10$ cm\,d$^{-1}$.}\sol \givenline{$q^\ast=71.5$ cm\,d$^{-1}$ ; $K_s=300$ and $10$ cm\,d$^{-1}$}
\par\noindent \[\phi_D=300/(300+71.5)=0.81\quad\text{\finans}\] for the permeable soil, and $\phi_D=10/(10+71.5)=0.12$ \finans for the tight one. The same stored water drains to four fifths in one soil and to one eighth in the other: this is the rate competition between drainage and the withdrawal by the crop and the atmosphere.

\section{The audit of the water balance and the compact law of evapotranspiration}
\textbf{The question.} Process-model outputs are used to support a saving of water, so they must be treated as closed accounts. The second study asks whether the water balance of every run closes, which term is missing when it does not, and whether the evapotranspiration of the ensemble can be summarised by a short law.

\par\noindent\textbf{The answer.} A pipeline first recovers the known water balance as a positive control, then audits the partition of the water among its fluxes: the seven fluxes obey linear invariants that the audit finds from the data, and what does not close is isolated as an unattributed term $R$, which is traced to canopy evaporation and to the rain of the year. The evapotranspiration of the ensemble is then summarised by a compact three-term law that is tested by leaving each experiment out in turn. The audit also shows that the compensation of irrigation, $\kappa$, is not the same in every soil, and it records the pattern of the 149 failed runs of the 10,190.

\par\noindent\textbf{How to read the figures.} They follow this order: the pipeline, the invariants, the unattributed term, the partition manifold, the compact law, the compensation and the fates of water by soil, and the failed runs.

\par\vspace{6pt}\noindent\hd{The pipeline for interpretable discovery}\par\nobreak
\fig{at/a22_16.jpg}{0.84}{Pipeline for interpretable discovery. The pipeline recovers the known water balance as a positive control, audits partitions, measures the dimension of the partition manifold and fits sparse laws with LOSO/LOEO validation against gradient boosting. Controls (dashed purple boxes) are shown at the bottom and outputs at the top; purple hexagons are processing blocks, the pink stippled box is the data store, peach boxes are outputs and dotted lines are flows of information.}\label{fig:a22_16}
\par\vspace{6pt}\noindent\hd{Linear invariants of the fluxes}\par\nobreak
\fig{at/a22_17.jpg}{0.84}{Linear invariants. (a) Singular values of the seven-flux matrix for the real data and for column-permuted data. (b) Coefficients of the smallest right singular vector (scaled so $c_{\mathit{ET}}=-1$) with bootstrap intervals. (c) Storage residual $\Delta S/W$ by experiment: E8 and E9 do not close. (d) Unattributed share of ET, $R/\mathit{ET}$, by experiment.}\label{fig:a22_17}
\par\vspace{6pt}\noindent\hd{The unattributed term}\par\nobreak
\fig{at/a22_18.jpg}{0.84}{The unattributed term $R$. (a) $R$ and canopy evaporation in E4 by irrigation method. (b) $R$ against annual rain in E9 (rainfed, light points), with the linear fit and the ten E0 baselines.}\label{fig:a22_18}
\par\vspace{6pt}\noindent\hd{The partition manifold}\par\nobreak
\fig{at/a22_19.jpg}{0.84}{Partition manifold and its sensitivity to zero handling. (a) Variance explained by the first three clr-PCA components for five zero-replacement floors. (b) Scores on the first two components (1\,mm floor) coloured by $\ln(W/\TAW)$. (c) Loadings of the five fates on the first two components.}\label{fig:a22_19}
\par\vspace{6pt}\noindent\hd{The compact law of evapotranspiration}\par\nobreak
\fig{at/a22_20.jpg}{0.84}{Compact ET law. (a) Leave-one-experiment-out predictions of the three-term law versus Daisy ET. (b) Gradient-boosting benchmark. (c) Fitted coefficients across the 17 LOEO folds. (d) LOEO MAPE by experiment. (e) The law at $I=200$\,mm for four soils (grey dashed: ET $=W$). (f) Drainage derived from closure, $W-\widehat{\mathit{ET}}$, versus Daisy drainage ($R^2=0.69$, RMSE 64\,mm; E8 and E9 excluded).}\label{fig:a22_20}
\par\vspace{6pt}\noindent\hd{Irrigation compensation}\par\nobreak
\fig{at/a22_21.jpg}{0.84}{Irrigation compensation $\kappa$ is not universal. (a) Per-soil $\kappa$ in E16 (paired $P=0.5\to1.5$) with the pooled estimate and bootstrap interval. (b) Slope of water input on rain in single-soil lattices (95\,\% intervals). (c) Mean $\kappa$ by factor level in E16.}\label{fig:a22_21}
\par\vspace{6pt}\noindent\hd{Water fates by soil}\par\nobreak
\fig{at/a22_22.jpg}{0.84}{Water fates by soil in the 15 near-equilibrium experiments (E8, E9 excluded). (a) Mean fate shares of water input $W$, soils ordered by decreasing $n$. (b) Number of valid runs per soil. (c) Median share of the unattributed term in ET.}\label{fig:a22_22}
\par\vspace{6pt}\noindent\hd{The pattern of the failed runs}\par\nobreak
\fig{at/a22_23.jpg}{0.84}{Pattern of the 149 failed runs. (a) Failed runs per soil, as a percentage of the 519 runs per soil (counts above bars). (b) Failed fraction in S02 against irrigation depth per event in E10. (c) Failed fraction in S02 against rain scale. (d) Annual drainage at 1\,m for valid and failed runs, normalised as densities; negative drainage of several hundred millimetres occurs only in failed runs.}\label{fig:a22_23}
\section{SWAFL Desktop: the engines and their verification}
\textbf{The question.} A model is useful only if its implementation can be trusted. The third study describes the application, SWAFL Desktop, and the checks that were made on it: that the design modules reproduce published worked examples, that the engines agree with Daisy on runs that were held out, and that the solvers agree with independent solvers.

\par\noindent\textbf{The answer.} The application has two engines, one for annual decisions and one that works on the daily state, and a set of design modules. Engine~1 is compared with Daisy on 1\,000 held-out runs, Engine~2 is checked for closure of its daily core and against Engine~1, the drip-lateral solver is compared with an independent converged solver on 400 random configurations, the surface-irrigation module reproduces the examples of the FAO Irrigation Paper 45, and a one-at-a-time sensitivity test shows what each control does. Daisy 7.1.14 runs end to end from the application, including a field of 5.85 ha taken from a satellite image to a schedule.

\par\noindent\textbf{How to read the figures.} They follow the checks: the overview, the verification against a worked example, the comparison with Daisy, the engines, the solver, the sensitivity, the published sources and the views of the application.

\par\vspace{6pt}\noindent\hd{Architecture of SWAFL Desktop}\par\nobreak
\fig{at/a22_24.jpg}{0.84}{Schematic overview of the main processes in SWAFL Desktop. Boxes indicate storages and solid lines flows; dashed black lines are flows out of the system, dotted lines are flows of information and dashed grey lines exchange of information between modules. Dashed purple boxes with a triangle are the drop-down controls and driving variables of the Properties panel, and the drivers of soil hydraulic properties (SHP). Dashed red lines and numbered circles mark the seven feedback loops.}\label{fig:a22_24}
\par\vspace{6pt}\noindent\hd{Verification against a textbook example}\par\nobreak
\fig{at/a22_25.jpg}{0.84}{Verification against Example 6.1. (a) Advance time $t_L$, (b) application efficiency $E_a$ (published versus SWAFL; grey line is 1:1), (c) relative differences of the 16 quantities (blue: advance and required times; green: efficiencies).}\label{fig:a22_25}
\par\vspace{6pt}\noindent\hd{Soil-water-atmosphere variables of the ten soils}\par\nobreak
\fig{at/a22_26.jpg}{0.84}{fig04 swa state}\label{fig:a22_26}
\par\vspace{6pt}\noindent\hd{Engine 1 on held-out runs}\par\nobreak
\fig{at/a22_27.jpg}{0.84}{Engine 1 versus Daisy for the 1\,000 held-out E12 runs. Orange: as shipped; blue: re-anchored. (a) ET, (b) irrigation, (c) drainage, (d) yield, (e) transpiration, (f) soil evaporation, (g) nitrogen leached. Grey lines are 1:1.}\label{fig:a22_27}
\par\vspace{6pt}\noindent\hd{Engine 2 and its closure}\par\nobreak
\fig{at/a22_28.jpg}{0.84}{Engine 2. (a) Closure residual of the unscaled daily core over 60 random states. (b) Mean annual irrigation of Engine 2 against Engine 1 for the 100 scenarios of the realistic grid (10 soils $\times$ 5 triggers $\times$ 2 methods, cereal), including the failures at zero. (c) Unscaled Engine 2 drainage against Engine 1 for the same scenarios.}\label{fig:a22_28}
\par\vspace{6pt}\noindent\hd{The drip-lateral solver}\par\nobreak
\fig{at/a22_29.jpg}{0.84}{Drip-lateral solver versus an independent converged solver (400 random configurations). (a) Relative difference in total head loss for the four-pass solver of version 1 and the converged solver, against the loss ratio $h_f/H_0$ (dashed: design limit 0.2; dotted: 0.6); (b) number of iterations of the converged solver; crosses mark configurations that did not converge.}\label{fig:a22_29}
\par\vspace{6pt}\noindent\hd{One-at-a-time sensitivity}\par\nobreak
\fig{at/a22_30.jpg}{0.84}{One-at-a-time sensitivity of SWAFL Desktop outputs. (a) Engine 1 annual irrigation and (b) drainage change for each control moved to the two settings listed under its name (blue: first setting; orange: second setting) from the reference state; nitrogen rate leaves both unchanged. (c) Irrigation and drainage by soil. (d) Drip emission uniformity against lateral length by lateral diameter and (e) by inlet head; dashed where $h_f/H_0>0.2$.}\label{fig:a22_30}
\par\vspace{6pt}\noindent\hd{Checks against published sources}\par\nobreak
\fig{at/a22_31.jpg}{0.84}{Checks against published sources. (a) Thirteen advance times of the FAO Irrigation Paper 45 border examples (published against SWAFL; grey line 1:1). (b) Christiansen coefficient of uniformity of the impact-sprinkler layer against spacing as a share of the wetted diameter, without wind and at 3\,m\,s$^{-1}$; dashed: 65\,\% spacing rule, dotted: CU of 85\,\%.}\label{fig:a22_31}
\par\vspace{6pt}\noindent\hd{Result views: the soil and the field}\par\nobreak
\fig{at/a22_32.jpg}{0.84}{fig views A}\label{fig:a22_32}
\par\vspace{6pt}\noindent\hd{Result views: irrigation timing and design}\par\nobreak
\fig{at/a22_33.jpg}{0.84}{fig views B}\label{fig:a22_33}
\par\vspace{6pt}\noindent\hd{Result views: mulch and the whole system}\par\nobreak
\fig{at/a22_34.jpg}{0.84}{fig views C}\label{fig:a22_34}
\section{The reach of the trigger, the fate of a millimetre and the dose}
\textbf{The question.} A sensor at a fixed depth starts an irrigation when the suction there crosses a trigger. How far above or below the sensor does the trigger really see, where does an irrigated millimetre end up, and how does the answer change with the size of the dose?

\par\noindent\textbf{The answer.} The \emph{reach} $h_c$ is the suction beyond which the sensor can no longer fire, and it is set by the soil through $\alpha$, $n$, $m$ and the connectivity $\ell$. The \emph{fate} of a millimetre is its division among drainage, evaporation, transpiration and storage, and for the drained share it grows with the dose as $\eta_D(d)=\eta_\infty\bigl(1-e^{-d/d_D}\bigr)$: small doses are kept by the soil, large doses leak. The trigger, the soil and the dose together decide whether an irrigation saves water or only moves it.

\par\noindent\textbf{How to read the figures.} They show the reach for ten soils, the fates as a function of the dose, and the way the sensor depth and the soil change the result.

\par\vspace{6pt}\noindent\hd{The two questions of the trigger}\par\nobreak
\fig{at/a22_35.jpg}{0.84}{The two questions. Q1 (right): a trigger $h_{\mathrm{tr}}$ fires only if the crop can dry the soil at the sensor to $h_{\mathrm{tr}}$; the reach $\hc$ (dashed) is 199\,cm in sand and $9\times10^{3}$\,cm in loam, so the usual triggers (grey) lie beyond the reach of sand. Q2 (left): the fate of an irrigation millimetre, averaged over the 744 irrigated runs of experiment E10 (arrow widths to scale).}\label{fig:a22_35}
\par\vspace{6pt}\noindent\hd{Retention and conductivity with the reach}\par\nobreak
\fig{at/a22_36.jpg}{0.84}{(a) Retention and (b) relative conductivity of the ten soils (the equation of the paper); circles mark the reach of the equation of the paper and the dashed line the fitted $\kappa=3.5\times10^{-10}$; grey: the range of triggers tested. (c) Reach against $n$ for the ten soils, with the equation of the paper for three values of $\alpha$.}\label{fig:a22_36}
\par\vspace{6pt}\noindent\hd{The reach of a suction trigger}\par\nobreak
\fig{at/a22_37.jpg}{0.84}{The reach of a suction trigger in experiment E2 (197 runs). (a) Irrigation against the trigger. (b) The same runs against the reach-scaled trigger $h_{\mathrm{tr}}/\hc$, irrigation divided by its soil maximum: every soil stops at $h_{\mathrm{tr}}/\hc\approx0.5$ to 1. (c) Predicted reach (the equation of the paper) against the observed interval (bars; arrows: only a lower bound); grey band: within a factor of 2. (d) Yield against the trigger: beyond the reach the crop is rainfed.}\label{fig:a22_37}
\par\vspace{6pt}\noindent\hd{The fate of an irrigation millimetre}\par\nobreak
\fig{at/a22_38.jpg}{0.84}{The fate of an irrigation millimetre, the equation of the paper. (a) By soil (E10, 744 irrigated runs; sand and loamy sand are never irrigated at the 300\,cm trigger of E10), ordered by $n$. (b) By method, on all soils (E4) and on coarse loam over all rates, doses and rooting depths (E15).}\label{fig:a22_38}
\par\vspace{6pt}\noindent\hd{The dose of irrigation}\par\nobreak
\fig{at/a22_39.jpg}{0.84}{The dose (E10, all ten windows pooled). (a) Drained share with the equation of the paper per soil. (b) Transpired share. (c) Evaporated share (soil and other evaporation).}\label{fig:a22_39}
\par\vspace{6pt}\noindent\hd{A trigger law independent of the soil}\par\nobreak
\fig{at/a22_40.jpg}{0.84}{A trigger law independent of the soil (E2, ten soils, twenty triggers). Relative irrigation and yield against the trigger, as suction and as root-zone depletion $f$; the curves of the ten soils collapse when depletion is used, and yield falls below 0.95 of its maximum beyond $f^{*}=0.45$. Panel d: the soil-specific suction that gives $f^{*}$, and the depletion reached by the 300\,cm trigger.}\label{fig:a22_40}
\par\vspace{6pt}\noindent\hd{Water fluxes as a Sankey diagram}\par\nobreak
\fig{at/a22_41.jpg}{0.84}{Mean annual water fluxes of E2 at the 300\,cm trigger on coarse loam (left) and coarse sand (right): inputs (rain, irrigation) and outputs (transpiration, soil evaporation, interception and other evaporation, drainage below 1\,m).}\label{fig:a22_41}
\par\vspace{6pt}\noindent\hd{How far management can move each outcome}\par\nobreak
\fig{at/a22_42.jpg}{0.84}{How far management can move each outcome, by soil. Controllability of yield, irrigation efficiency, drainage and nitrate-leaching fractions by seven management levers (E1 to E7) across the ten soils, and its rank correlation with the plant-available water of the soil.}\label{fig:a22_42}
\exbox{The dose of an irrigation event}{The drained share of an irrigation dose $d$ follows $\eta_D=\eta_\infty(1-e^{-d/d_D})$ with a plateau $\eta_\infty=0.15$ and a dose scale $d_D=4.1$ mm. Find the drained share for doses of 1 mm and 10 mm.}\sol \givenline{$\eta_\infty=0.15$ ; $d_D=4.1$ mm ; $d=1,\ 10$ mm}
\par\noindent For $d=1$ mm, \[\eta_D=0.15\,(1-e^{-0.244})\ti{i}\]
Substituting the values in equation \rfi{i},
\[\eta_D=0.15\times0.217=0.032\]
for $d=10$ mm, \[\eta_D=0.15\,(1-e^{-2.44})\ti{ii}\]
Substituting the values in equation \rfi{ii},
\[\eta_D=0.15\times0.913=0.137\quad\text{\finans}\]
A dose of 1 mm drains $3$ per cent and one of 10 mm drains $14$ per cent, almost the plateau: small doses are lost by evaporation, large doses by drainage, which is the result of the dose figure.

\section{The workbench, the hand-over, the field and the plants}
\textbf{The question.} How does a user go from a soil, a weather file and a crop to a decision, and how is the workbench that does it checked? The last two studies describe the workbench of the application (the modules that set up and run Daisy), the hand-over of a run to the decision layers, and the tests of the numerical choices that the workbench makes.

\par\noindent\textbf{The answer.} The workbench builds a Daisy run from forms, runs it, and passes the log to the layers, so that the same water balance feeds the schedule, the design and the audit. The compact designs that it uses to save time are tested against the marginal curves of the full study; the soil-column solver is tested for the effect of the node spacing; the denitrification function is checked against 792 stored Daisy outputs. The crop library and the plants close the chapter: the grain mass of seven genotypes under a drought pulse, and the daily crop-coefficient curves of the library.

\par\noindent\textbf{How to read the figures.} They follow the workbench module by module, then the numerical tests, then the architecture of the application, and last the crops and the plants.

\par\vspace{6pt}\noindent\hd{Modules of the Daisy workbench}\par\nobreak
\fig{at/a22_43.jpg}{0.84}{Modules of the Daisy workbench and the hand-over. The chain runs from the planner and the decision base through the experiments, Daisy and the parsers to the hand-over; the modules below the dotted line are used before Daisy is started.}\label{fig:a22_43}
\par\vspace{6pt}\noindent\hd{The process map of Daisy}\par\nobreak
\fig{at/a22_44.jpg}{0.84}{The process map of Daisy as drawn by the workbench. Each process is a button that opens the figures it produces from a run.}\label{fig:a22_44}
\par\vspace{6pt}\noindent\hd{Water content of a soil column}\par\nobreak
\fig{at/a22_45.jpg}{0.84}{Water content of a soil column through time, computed by the Richards solver of the workbench.}\label{fig:a22_45}
\par\vspace{6pt}\noindent\hd{Fidelity of the compact designs}\par\nobreak
\fig{at/a22_46.jpg}{0.84}{(a) Error of the marginal curves rebuilt from the levels of the compact designs (median over five responses, per experiment and factor). (b) Daisy run time of the study designs and of the compact designs, by experiment (logarithmic scale).}\label{fig:a22_46}
\par\vspace{6pt}\noindent\hd{The Richards solver and its node spacing}\par\nobreak
\fig{at/a22_47.jpg}{0.84}{Soil column. (a) Drainage below 1\,m with the node spacing; dashed lines are the reference. (b) Final storage minus the reference. (c) Infiltration of 200\,mm of rain into a ponded silty clay loam (left axis) and run time (right axis), showing the lack of convergence.}\label{fig:a22_47}
\par\vspace{6pt}\noindent\hd{The denitrification function}\par\nobreak
\fig{at/a22_48.jpg}{0.84}{Denitrification function against the 792 stored Daisy outputs. (a) N$_2$ loss; the 1:1 line is black. (b) Histogram of the relative error of the N$_2$O loss before and after the correction.}\label{fig:a22_48}
\par\vspace{6pt}\noindent\hd{The bridge from a Daisy run to a decision}\par\nobreak
\fig{at/a22_49.jpg}{0.84}{The bridge from a Daisy run to a decision. Heavy arrows are the path of a run, dashed arrows are outputs.}\label{fig:a22_49}
\par\vspace{6pt}\noindent\hd{The economics of the study}\par\nobreak
\fig{at/a22_50.jpg}{0.84}{Daisy wall-clock time of the 17 experiments of the study, and the time of the compact designs of the workbench compared with it.}\label{fig:a22_50}
\par\vspace{6pt}\noindent\hd{Architecture of SWAFL Desktop 1.1.0}\par\nobreak
\fig{at/a22_51.jpg}{0.84}{Architecture of SWAFL Desktop 1.1.0. The field tools call the engines and write to the project; the project carries the state of the field.}\label{fig:a22_51}
\par\vspace{6pt}\noindent\hd{Drought and the grain mass of seven genotypes}\par\nobreak
\fig{at/a22_52.jpg}{0.84}{Grain dry mass at maturity of seven genotypes, well-watered and under a drought pulse, in a warm and in a hot, dry season.}\label{fig:a22_52}
\par\vspace{6pt}\noindent\hd{Crop coefficient curves of the library}\par\nobreak
\fig{at/a22_53.jpg}{0.84}{Daily crop-coefficient curves of the library at the default sowing date: five annual crops and three perennial crops.}\label{fig:a22_53}
\IfFileExists{extra/plus22.tex}{\banner{The methods of the research: design of experiments, validation and uncertainty}\label{ch:res-deep}
\section{Designing a simulation experiment}\label{sec:res-doe}
A model experiment is designed like a field trial. The \kw{factors} are the inputs that are varied (trigger, dose, mulch, rain), the \kw{levels} are the values chosen for each, and the \kw{responses} are the outputs that are compared. Three designs were used in this study.
\begin{enumerate}[leftmargin=2em,itemsep=2pt,label={\bfseries\color{khmag}\arabic*.}]
\item \hd{One factor at a time.} Each factor is varied over many levels with the others fixed (E1 doses, E2 triggers). Cheap and easy to read, but blind to interactions.
\item \hd{Full factorial.} Every combination of every level: $k$ factors at two levels need $2^k$ runs per soil (E16: $2^6=64$). It estimates all main effects and interactions.
\item \hd{Latin hypercube.} For many factors at many levels, $N$ runs are drawn so that each factor is sampled once in each of $N$ equal strata. The design covers the space evenly with far fewer runs than a factorial; SWAFL uses it with a fixed seed to keep a Daisy experiment within its run budget.
\end{enumerate}

\begin{kfig}{Two designs for two factors. Left: a full factorial of $3\times3$ levels (9 runs). Right: a Latin hypercube of 6 runs, one in each row and each column of the strata.}
\draw[ln] (0,0) rectangle (3,3);\foreach \x in {0.5,1.5,2.5}{\foreach \y in {0.5,1.5,2.5}{\fill[khmag] (\x,\y) circle (0.09);}}
\node[lab,below] at (1.5,-0.1) {factor 1};\node[lab,rotate=90] at (-0.3,1.5) {factor 2};\node[lab,above] at (1.5,3.05) {full factorial};
\begin{scope}[xshift=4.5cm]
\draw[ln] (0,0) rectangle (3,3);\foreach \i in {0.5,1,...,2.5}{\draw[thin ln,dashed] (\i,0)--(\i,3);\draw[thin ln,dashed] (0,\i)--(3,\i);}
\foreach \x/\y in {0.3/1.7,0.8/0.2,1.2/2.8,1.7/1.1,2.3/0.6,2.7/2.2}{\fill[khmag] (\x,\y) circle (0.09);}
\node[lab,below] at (1.5,-0.1) {factor 1};\node[lab,above] at (1.5,3.05) {Latin hypercube};
\end{scope}
\end{kfig}

\exbox{Counting the runs of a design}{(a) How many runs does a two-level full factorial of six factors need for ten soils? (b) If thirteen extra edge cases are run per soil, what is the total? (c) How many runs would a three-level factorial of the same six factors need for one soil?}\sol
\par\noindent (a) $2^6\times10=640$ runs.\finans\quad (b) $640+13\times10=770$ runs.\finans\quad (c) $3^6=729$ runs for one soil.\finans
\par\noindent The jump from 64 to 729 runs per soil when one level is added is why Latin hypercubes are used once more than a few factors are varied.

\section{Validation with each soil left out}\label{sec:res-loso}
A law fitted to ten soils will always fit them well; the question is whether it predicts a soil that it has not seen. In \kw{leave-one-soil-out} (LOSO) validation the law is fitted ten times, each time without one soil, and the left-out soil is predicted. The error measures of Section~\ref{sec:at-rcl} are then computed on the ten predictions. A law that does well in-sample and badly in LOSO is \kw{overfitted}: it has learnt the soils, not the process.

\begin{kfig}{Leave-one-soil-out validation: in each of ten rounds nine soils (blue) are used to fit the law and the tenth (magenta) is predicted.}
\foreach \r in {0,...,4}{\foreach \c in {0,...,9}{\fill[khcyan!50] (\c*0.62,-\r*0.5) rectangle ++(0.5,0.36);}\fill[khmag] (\r*0.62,-\r*0.5) rectangle ++(0.5,0.36);\node[lab,left] at (-0.1,{-\r*0.5+0.18}) {round \pgfmathparse{int(\r+1)}\pgfmathresult};}
\node[lab] at (3.0,-2.55) {$\vdots$};\foreach \c in {0,...,9}{\node[lab,above] at ({\c*0.62+0.25},0.4) {S\pgfmathparse{int(\c+1)}\pgfmathresult};}
\end{kfig}

\section{Uncertainty by the bootstrap}\label{sec:res-boot}
A fitted constant, such as $q^*$ or $\kappa$, is uncertain because the soils are a sample. The \kw{bootstrap} draws, many times, a new set of ten soils from the ten with replacement (so a soil may appear twice and another not at all), refits the constant on each set and takes the 2.5 and 97.5 percentiles of the refitted values as the 95 per cent interval. For $q^*$, 4999 resamples gave 58 to 105 cm\,d$^{-1}$ around the estimate of 71.5. The interval is not symmetric because the law is non-linear in $q^*$.

\exbox{How sure is the constant?}{The bootstrap interval of $q^*$ is 58.4 to 105.2 cm\,d$^{-1}$. Find the range of the predicted $\phi_D$ of the loam ($K_s=24.96$ cm\,d$^{-1}$) over this interval, and compare it with the change of $\phi_D$ between the loam and the sandy loam ($K_s=106.1$, $\phi_D=0.60$).}\sol
\par\noindent At $q^*=58.4$: $\phi_D=24.96/83.36=0.299$; at $q^*=105.2$: $\phi_D=24.96/130.16=0.192$.\finans
\par\noindent The uncertainty of the constant moves the loam by about 0.1, while the difference between the two soils is 0.34: the ranking of the soils is robust even though the constant is uncertain.
}{}
\section{Summary diagrams of the studies}
Each of the studies ends in a one-picture summary, a \emph{graphical abstract}, that gives its question, its method and its main result in a single diagram. They are collected here as a map of the chapter: a reader who has followed the preceding sections can use them to recall, in a minute, what each study found and where its result is used in the application.

\par\vspace{6pt}\noindent\hd{Graphical abstract of the first paper}\par\nobreak
\fig{at/a22_54.jpg}{0.84}{Graphical abstract of the first paper: the antecedent storage of the root zone and the three places where it can go.}\label{fig:a22_54}
\par\vspace{6pt}\noindent\hd{Graphical abstract of the second paper}\par\nobreak
\fig{at/a22_55.jpg}{0.84}{Graphical abstract of the second paper: the audit of the water balance and the term that no box of the model carries.}\label{fig:a22_55}
\par\vspace{6pt}\noindent\hd{Graphical abstract of the third paper}\par\nobreak
\fig{at/a22_56.jpg}{0.84}{Graphical abstract of the third paper: the engines of SWAFL Desktop, from the soil to the decision.}\label{fig:a22_56}
\begin{itemize}\item The six papers report one study: the memory of soil water, the audit of the model, the engines of SWAFL, the reach of the trigger, the workbench and the field tools.\item A figure of a paper is read with its caption: the panels (a), (b), (c) are separate statements.\item Every figure can be redrawn from the scripts and the tables of the repository.\end{itemize}
\backhead{EXERCISES}{khred}\begin{exlist}\item Name the three places where the antecedent storage of the root zone can go. \exans{drainage, soil evaporation, transpiration (and other evaporation)}\item In which figure do the ten soils collapse onto one curve? \exans{depletion trigger law}\item Which figure compares the engines of SWAFL with Daisy on held-out runs? \exans{Engine 1 held-out, Paper 3}\end{exlist}
\IfFileExists{extra/more22x.tex}{\backhead{FURTHER EXERCISES}{khred}
\begin{exlist}
\item How many runs does a two-level full factorial of eight factors need per soil? A Latin hypercube of 50 runs covers each factor in how many strata?\exans{256; 50}
\item A bootstrap of 4999 resamples gives the 2.5 and 97.5 percentiles of a constant. Which resamples (by rank) are these?\exans{the 125th and the 4875th}
\item In leave-one-soil-out validation of ten soils, how many fits are made, and on how many soils is each fitted?\exans{10; 9}
\item Two laws have LOSO RMSE 0.058 and 0.111. By what percentage is the first better?\exans{48 per cent}
\item A constant has the interval 58 to 105 around 71.5. Is the interval symmetric? Find the distances below and above.\exans{no; 13.5 and 33.5}
\end{exlist}
\backhead{ESSAY QUESTIONS}{khred}
{\small\itshape Answer each question in about one page, with equations, a sketch and at least one number from this chapter.}\par\vspace{3pt}
\begin{exlist}
\item Describe the design of the simulation experiments of the study and the reasons for each design.
\item Discuss leave-one-out validation and why it is the main test of the laws of this book.
\item Explain the bootstrap and what its intervals mean.
\item Discuss the verification of SWAFL against textbook examples and held-out runs.
\item Explain how the studies of the project fit together into one argument.
\item Choose one study and write a critical review of its methods and conclusions.
\end{exlist}
}{}
\dykband{DO YOU KNOW ?}\begin{exlist}\item What is the unattributed term of the water balance?\item What does the reach of a trigger mean?\item How is the fate of an irrigation millimetre defined?\item Why is the dose of irrigation important?\end{exlist}
\IfFileExists{extra/more22s.tex}{\backhead{OBJECTIVE TYPE QUESTIONS}{khred}
\begin{exlist}
\item A two-level full factorial of six factors has\opt{12}{36}{64}{720} runs
\item In leave-one-soil-out validation of ten soils each fit uses\opt{1}{5}{9}{10} soils
\item A bootstrap draws soils\opt{without replacement}{with replacement}{in order}{once each}
\item A Latin hypercube samples each factor\opt{at two levels}{once in each stratum}{at random only}{at its mean}
\item An overfitted law has\opt{a low LOSO error}{a low in-sample error and a high LOSO error}{no parameters}{a high in-sample error}
\end{exlist}
\objkey{1.\ (c),\quad 2.\ (c),\quad 3.\ (b),\quad 4.\ (b),\quad 5.\ (b)}
\dykband{EXERCISES IN SWAFL}
{\small Open SWAFL (desktop or online) with the benchmark project. The chapter of the \href{https://github.com/halilunaziru73-creator/Coupled-Water-Fluxes-Pore-Scale-Dynamics-Precision-Irrigation-Mulching-Feedback-Loops/raw/main/Software/v1.1.0/SWAFL_Desktop/docs/manual_v1.1.0/SWAFL_Manual_v1.1.0.pdf}{SWAFL manual} that explains each task is given in brackets; record your values in a table and compare them with what is expected (\emph{Check}).}\par\vspace{3pt}
\begin{exlist}
\item {\color{khblue}[Manual, Chapter 28]} Design a Daisy experiment with a Latin hypercube of 20 runs and inspect the plan.\par\noindent{\small\itshape Check: each factor is sampled in 20 strata.}
\item {\color{khblue}[Manual, Chapter 29]} Open the best combinations and choose one for your field.\par\noindent{\small\itshape Check: the plan of the combination is shown.}
\item {\color{khblue}[Manual, Chapter 30]} Read the interactions of an experiment and identify the strongest.\par\noindent{\small\itshape Check: usually between water and nitrogen.}
\item {\color{khblue}[Manual, Chapter 18]} Use the search for the most efficient inputs and compare with the best combination of Daisy.\par\noindent{\small\itshape Check: the two agree in direction.}
\end{exlist}
}{}
\swaflvideo{22}
